\chapter{Introdução}

\section{Contextualização e delimitação do tema}

A incorporação de recursos de inteligência artificial às rotinas escolares e sociais dos estudantes aproxima os jovens de respostas prontas que podem ser usadas para concluir tarefas sem que se compreenda o conteúdo envolvido. Esse movimento se intensifica nas redes sociais e em plataformas de vídeo, onde circulam demonstrações de ferramentas que realizam atividades acadêmicas e prometem resultados rápidos.

Assistir a essas demonstrações, porém, não garante que o estudante aprenda o conteúdo nem que compreenda como a inteligência artificial opera, quais dados utilizou e por qual processo chegou ao resultado. Sem esse entendimento, o estudante pode ter dificuldade para analisar criticamente a resposta recebida, e o currículo escolar pode oferecer pouco espaço para explicar como um sistema de inteligência artificial chega a seus resultados. A Resolução CNE/CEB nº 1/2022 define diretrizes para a Computação na Educação Básica como complemento à Base Nacional Comum Curricular e atribui aos sistemas de ensino a organização da implantação curricular \cite{brasil2022computacao}. Em 2026, o Ministério da Educação publicou o documento orientador \textit{Inteligência artificial na Educação Básica}, que reforça o papel do professor e defende que a IA seja tratada como recurso de tutoria e apoio à aprendizagem, e não apenas como geradora de respostas prontas \cite{brasil2026iaeducacaobasica}.

A educação em inteligência artificial tem sido debatida internacionalmente como parte da formação cidadã, pois os estudantes precisam compreender, conforme sua escolaridade, as possibilidades e os limites de sistemas inteligentes \cite{touretzky2019envisioning}. No Ensino Médio, essa necessidade se coloca de modo particular: usar muito os recursos digitais não é o mesmo que ter letramento computacional. Muitos estudantes usam celulares e aplicativos com desenvoltura, mas nem sempre tiveram oportunidade de interpretar com autonomia uma tabela, distinguir correlação de previsão ou perceber por que um resultado muda quando a regra é alterada. Outros ainda enfrentam restrições de acesso a computador e internet fora da escola.

A discussão sobre o uso responsável deve ser acompanhada pelo ensino dos conceitos que explicam como a inteligência artificial gera suas respostas. Se a conversa se limita a liberar ou proibir uma ferramenta, o estudante permanece na posição de usuário. Compreender como uma inteligência artificial é treinada pode ajudá-lo a entender como ela produz respostas e a utilizá-la de modo favorável ao seu desenvolvimento cognitivo. O ensino de Computação permite deslocar a conversa para o que alimenta um modelo e por que uma mudança nos dados altera sua resposta, e este projeto focaliza o ensino introdutório de \textit{machine learning} para discuti-lo a partir da relação entre dados e modelos computacionais.

O recorte não pretende exigir programação prévia nem formar especialistas em algoritmos. Em oito encontros, estudantes do Ensino Médio observarão dados, formularão hipóteses e explicarão o efeito de modificar parâmetros, com a Matemática como linguagem de interpretação, e não como barreira inicial para o estudo do problema.

\section{Problema educacional e justificativa}

Na rotina escolar, alguns estudantes recorrem à inteligência artificial para obter resultados imediatos, sem necessariamente se apropriar do conhecimento envolvido. Quando esse uso se limita à obtenção de respostas prontas, o estudante não necessariamente compreende como a ferramenta funciona. Nesse caso, o ensino dos conceitos básicos de \textit{machine learning} pode ajudá-lo a conhecer o processo por trás das respostas que recebe. Outro desafio é oferecer aos estudantes oportunidades de trabalhar com dados.

Em geral, tabelas e gráficos apresentados no Ensino Médio chegam prontos, inseridos em questões. Nessas situações, a leitura dos valores pode ocorrer sem uma discussão sobre a origem dos dados ou sobre o que mudaria no resultado se um valor ou uma condição fosse diferente.

Compreender um modelo de \textit{machine learning} exige algum contato prévio com dados, além de noções básicas de estatística e probabilidade. Já existem trabalhos voltados ao ensino do tema na Educação Básica. Em um mapeamento de 30 unidades didáticas, \citeonline{marques2020teaching} observaram que várias delas cobrem apenas os processos mais acessíveis, como a gestão de dados, ou apresentam o aprendizado e o teste dos modelos de forma abstrata, deixando parte do funcionamento interno oculto. Isso sugere que, nessas unidades, o contato com dados pode se restringir à gestão, sem mostrar como os dados e os parâmetros alteram o resultado do modelo. No Ensino Médio, a revisão de \citeonline{martins2023findings} indica que os estudantes conseguem compreender e aplicar conceitos, algoritmos e tarefas básicas de \textit{machine learning}. A mesma revisão aponta a necessidade de avaliações mais rigorosas sobre como o tema é ensinado. A presença da inteligência artificial na escola também exige situações em que os estudantes relacionem os dados à finalidade da tarefa, à regra, ao parâmetro e ao resultado do modelo.

Quando os dados chegam prontos e as relações entre eles não são trabalhadas, o estudante tende a aprender os conceitos de \textit{machine learning} apenas pelo nome. Termos como treinamento, modelo, parâmetro e previsão podem ser repetidos sem que se saiba como se ligam entre si. Um estudante pode dizer que um modelo foi treinado sem saber explicar o que muda quando um parâmetro é alterado, e pode reconhecer uma reta no gráfico sem relacioná-la a uma previsão. O problema educacional está na passagem entre nomear um conceito e utilizá-lo para explicar uma situação.

Há ainda uma diferença importante entre apresentar a inteligência artificial como tema de cultura digital e ensiná-la como conteúdo de Computação. A discussão cultural é indispensável para tratar de riscos, vieses, autoria e impactos sociais. No ensino de Computação, o estudante também precisa investigar, na prática, uma representação de dados e uma regra computacional. A proposta não reduz questões éticas a uma equação ou a um gráfico. Considera-se que a discussão sobre limites passa a ter outra base quando o estudante reconhece que modelos resultam de decisões humanas de modelagem.

Romper com a recepção passiva de resultados prontos exige oportunidades para prever, alterar dados ou parâmetros, observar o efeito produzido e explicar por que ele surgiu. Esse processo começa antes do modelo, na origem dos dados. Um estudante que monta uma tabela com medidas da própria turma sabe de onde vêm os números e o que cada coluna representa. Essa experiência torna a leitura do gráfico menos mecânica. Em seguida, ver uma reta se deslocar quando os dados mudam, ou a mudança de categoria quando o limiar é alterado, produz evidências visíveis. Ao explicar o que observou, o estudante pode relacionar o efeito visual à alteração que realizou. Essa organização se aproxima das abordagens práticas, baseadas em problemas ou projetos, que \citeonline{martins2023findings} identificaram como bem-sucedidas no Ensino Médio. Compreender de onde vêm os dados e o que muda com um parâmetro é o que permite ao estudante questionar o resultado que recebe, e não apenas aceitá-lo.

O produto educacional prevê uma plataforma interativa em português brasileiro, articulada a uma sequência didática para estudantes do Ensino Médio. A plataforma apresenta cinco módulos introdutórios: regressão linear, regressão logística, \textit{k}-vizinhos mais próximos, árvore de decisão e \textit{k}-means. Os modelos serão tratados como conteúdos centrais, sem receber necessariamente a mesma formalização matemática. A regressão linear permite explorar tendência, previsão e erro; a regressão logística introduz probabilidade e limiar. O \textit{k}-NN coloca em evidência distância e votação; a árvore de decisão trabalha com regras e divisões; o \textit{k}-means permite discutir grupos e dados sem rótulos prévios.

O uso da plataforma, por si só, não será tomado como evidência de aprendizagem. Ferramentas visuais podem apoiar a compreensão, mas seu potencial depende de como as representações são desenhadas e das tarefas cognitivas pedidas aos estudantes \cite{ainsworth2006deft}. Por isso, cada módulo relacionará uma situação-problema, a origem dos dados, uma hipótese inicial, uma manipulação controlada e uma explicação posterior. A sequência segue o movimento de prever, experimentar, observar e explicar. A devolutiva do experimento participa da mediação docente, aproximando-se da avaliação formativa discutida por \citeonline{black1998assessment}. Serão oito encontros de cerca de 50 minutos. No segundo, a turma poderá organizar uma tabela simples antes de estudar regressão linear; nos demais módulos, trabalhará com cenários sintéticos fornecidos pelo site. Um estudante que modifica o limiar de uma classificação precisa discutir por que a mudança ocorreu. Sem essa discussão, a manipulação do controle pode ficar desvinculada do conceito que se pretende ensinar.

A justificativa profissional está na necessidade de um recurso que possa ser aplicado e analisado por um professor da Educação Básica, em turmas heterogêneas, com tempo limitado e infraestrutura que varia entre escolas. Para um programa profissional como o ProfComp, a articulação entre problema escolar, conteúdo de Computação, intervenção docente e produto utilizável é parte central da relevância do estudo. A plataforma não será avaliada como software genérico, nem a pesquisa pretende demonstrar que qualquer interface interativa resolve dificuldades de aprendizagem.

Espera-se contribuir para uma discussão que não trate a inteligência artificial como conteúdo reservado a cursos técnicos, nem a simplifique a ponto de ocultar como um modelo funciona. A contribuição é localizada e dependerá das condições de aplicação. Documentar essas condições pode ajudar outros docentes a avaliar se uma proposta semelhante serve ao seu contexto, sem antecipar generalizações que a pesquisa não poderá sustentar.

\section{Pergunta de pesquisa}

Que mudanças e permanências podem ser identificadas nas explicações de estudantes do Ensino Médio sobre fundamentos introdutórios de \textit{machine learning} ao longo de uma sequência didática apoiada por uma plataforma interativa?

A pergunta dirige a análise para as explicações, as interpretações de gráficos e tabelas e as escolhas justificadas pelos estudantes em diferentes momentos da sequência. O estudo não medirá um efeito isolado da plataforma. Examinará mudanças, permanências e dificuldades na forma como os participantes relacionam dados, finalidade da tarefa, regra, parâmetro e resultado. As produções serão elaboradas em um percurso que inclui a plataforma, a mediação do professor, o trabalho em dupla e as condições da turma.

